\section{Introduction} \label{sec:intro}
\begin{figure}[t]
  \centering
  \begin{minipage}[b]{0.8\linewidth}
    \includegraphics[width=\linewidth]{fig/fig_corruption_strip.pdf}
  \end{minipage}\\[6pt]
  \begin{minipage}[b]{0.8\linewidth}
    \includegraphics[width=\linewidth]{fig/fig_balanced_asymmetry_intro_delta.pdf}
  \end{minipage}
  \caption{Directional asymmetry in corrupted-gallery retrieval. \textbf{Top:} ToyCorruption severity ladder: one identity at $\sigma{=}0$ (clean) through $\sigma{=}4$ (extreme), produced by compositing Gaussian blur, resolution downscaling, gamma shift, sensor noise, and JPEG recompression following ImageNet-C~\cite{hendrycks2019benchmarking}. \textbf{Bottom:} Per-severity directional mAP gap $\Delta[k]{=}\mathrm{mAP}_A[k]{-}\mathrm{mAP}_B[k]$ under matched query and gallery counts (50 each, one correct match per query). Direction~A (clean query, corrupted gallery) is easier than Direction~B at $k{=}1$--$4$, with the gap peaking near $+0.20$ at $k{=}3$; this asymmetry motivates direction-aware scoring.}
  \label{fig:balanced_asymmetry_intro}
\end{figure}
Domain-generalizable person re-identification (DG-ReID) learns representations from multiple labeled source domains that can generalize to unseen target domains without additional fine-tuning \cite{cho2024generalizable}. Practical deployments mix resolutions, viewpoints, and camera parameters; much of the literature therefore stresses \textit{domain-invariant} features. This is often done by separating identity-discriminative cues from domain-specific appearance via normalization \cite{pan2018two,jin2020style}, meta-learned domain-invariant mapping \cite{song2019generalizable,choi2021meta}, or disentanglement-style factorizations \cite{zheng2019joint} (see Sec.~\ref{sec:relatedwork}). These approaches typically still rank with symmetric distances on Euclidean embeddings, including cosine- or spherical-normalized views.

We focus on a complementary modeling choice: \textbf{direction-aware retrieval} when probe-to-gallery matching is not symmetric in deployment (e.g., aerial-to-ground or cross-modality RGB-to-infrared). Symmetric distances impose $d(\mathbf{x},\mathbf{y})=d(\mathbf{y},\mathbf{x})$, which may be an acceptable abstraction for bidirectional search but can discard ordered degradation patterns across domains. It treats probe→gallery and gallery→probe as interchangeable when the operative query direction is fixed at test. Our goal is not to dismiss invariance-based DG-ReID, but to ask when an asymmetric score, viewed through a Randers/Finsler lens \cite{bao2000introduction}, better matches that directional geometry than a symmetric surrogate.

Building on Balancing Alignment and Uniformity (BAU) \cite{cho2024generalizable}, we introduce (i) a structured embedding with an identity factor and an orthogonalized drift factor, and (ii) an asymmetric distance on those factors with a bounded learnable weight that interpolates toward the Euclidean triplet regime. This keeps the representation geometry transparent: features remain Euclidean while asymmetry enters exclusively through the pairwise scoring function. We first study explicit disentanglement with nuisance-aligned auxiliary objectives on the drift factor. In our training stack, such objectives can interfere with domain-invariant terms: a failure mode analyzed in Sec.~\ref{subsec:interference_analysis} and a cautionary datapoint for “hard” disentanglement. That analysis motivates retaining a single continuous feature map while letting asymmetry enter through the scoring function and a stable triplet parameterization. 

% Our main contributions are summarized as follows: (i) we foreground asymmetric probe–gallery matching alongside domain invariance as a complementary design axis; (ii) we instantiate a Randers/Finsler-style asymmetric score on a split identity/drift representation with norm-constrained drift that recovers Euclidean geometry at initialization, atop BAU-style alignment and uniformity training; (iii) we analyze how strong nuisance alignment in the drift factor can clash with domain-invariant losses; (iv) we empirically stress-test the full scoring stack under fixed multi-source protocols, including the cross-camera drift coherence loss $\mathcal{L}_{\mathrm{dcc}}$ (Sec.~\ref{subsec:drift_aux_losses}); (v) we report an auxiliary-loss sweep that selects a primary training recipe and documents which multi-target and cross-view evaluations are still outstanding (Sec.~\ref{sec:experimentalResults}).

Our main contributions are summarized as follows: (i) we foreground asymmetric probe–gallery matching alongside domain invariance as a complementary design axis; (ii) we instantiate a Randers/Finsler-style asymmetric score on a split identity/drift representation with norm-constrained drift that recovers Euclidean geometry at initialization, atop BAU-style alignment and uniformity training; (iii) we analyze how strong nuisance alignment in the drift factor can clash with domain-invariant losses; (iv) we empirically stress-test the full scoring stack under fixed multi-source protocols, including the cross-camera drift coherence loss $\mathcal{L}_{\mathrm{dcc}}$ (Sec.~\ref{subsec:drift_aux_losses}); (v) we report an auxiliary-loss sweep that selects a primary training recipe and documents which multi-target and cross-view evaluations are still outstanding (Sec.~\ref{sec:experimentalResults}); (vi) we validate the asymmetric scoring mechanism in a controlled synthetic setting: a pairwise monotonicity loss encodes corruption severity in a learned scalar; the Randers correction produces the predicted directional distance gap; and balanced bidirectional evaluation at matched counts confirms that clean-probe retrieval is consistently easier than corrupted-probe retrieval (Sec.~\ref{subsec:controlled_corruption}).


% -- Extras / Random Thoughts / Previous Drafts -- 


